\chapter{运算的优先级}\label{ch:g7:priorities}

$3 + 4 \times 5$ 等于多少？若从左到右计算，得 $35$；
若先做乘法，则得 $23$。数学要求只有一个答案，
因此计算需要一套“交通规则”——即优先级。
本章给出这些规则，并引入第一个代数恒等式，
即分配律。

\section{优先级规则}

\begin{theorem}[优先级规则]\label{thm:g7:priorities:rules}
在不含括号的运算中：
\begin{enumerate}
  \item 乘法和除法\emph{先于}
        加法和减法进行；
  \item 同级运算按从左到右的顺序进行。
\end{enumerate}
括号内的运算优先于其他一切：先算括号内的部分，
并从最内层括号开始。
\end{theorem}

\begin{example}\label{ex:g7:priorities:basic}
每行只做一步，并在心中标出当前计算的部分：
\begin{align*}
3 + 4 \times 5 &= 3 + 20 = 23
&&\text{（先做乘法）}\\
20 - 8 + 2 &= 12 + 2 = 14
&&\text{（同级运算：从左到右）}\\
18 \div 3 \times 2 &= 6 \times 2 = 12
&&\text{（同级运算：从左到右）}\\
(3 + 4) \times 5 &= 7 \times 5 = 35
&&\text{（先算括号）。}
\end{align*}
注意易错之处：$20 - 8 + 2$ \emph{并不等于} $20 - 10$，
而 $18 \div 3 \times 2$ \emph{并不等于} $18 \div 6$。
\end{example}

\begin{example}[多层括号]\label{ex:g7:priorities:nested}
\begin{align*}
50 - 2 \times \bigl(3 + (10 - 4)\bigr)
&= 50 - 2 \times (3 + 6) && \text{（最内层括号）}\\
&= 50 - 2 \times 9 && \text{（剩下的括号）}\\
&= 50 - 18 && \text{（乘法）}\\
&= 32 .
\end{align*}
\end{example}

\begin{remark}[分数线会隐藏括号]
\omterm{def:g6:fractions:def}{分数}线将\omterm{def:g4:fractions:def}{分子}与\omterm{def:g4:fractions:def}{分母}各自视为一组，
仿佛各自带有看不见的括号：
\[
\frac{12 + 8}{4} = \frac{20}{4} = 5,
\qquad\text{而非}\quad 12 + \frac{8}{4} = 14 .
\]
\end{remark}

\begin{method}[安全计算]\label{met:g7:priorities:safe}
\begin{enumerate}
  \item 找出应先进行的运算（依次为最内层括号、
        $\times$ 与 $\div$、$+$ 与 $-$）；
  \item \emph{只计算这一步}，并将整个式子改写为
        结果放在原位后的形式；
  \item 重复上述步骤，直到只剩一个数——每行只做一步运算，
        不跳步。
\end{enumerate}
\end{method}

\section{分配律}

\begin{theorem}[分配律]\label{thm:g7:priorities:distributivity}
对所有数 $k$、$a$、$b$：
\[
k \times (a + b) = k \times a + k \times b,
\qquad
k \times (a - b) = k \times a - k \times b .
\]
\end{theorem}

\begin{proof}[图形证明]
计算宽为 $k$ 的长方形的\omterm{def:g6:measure:area}{面积}，该长方形被分为两个
较小的长方形，长分别为 $a$ 和 $b$：整体\omterm{def:g6:measure:area}{面积}为 $k \times (a+b)$；
分两部分计算时为 $k \times a + k \times b$。同一个长方形，
\omterm{def:g6:measure:area}{面积}自然相同。
\end{proof}

\begin{omfigure}
\begin{tikzpicture}[scale=0.75]
  \fill[omDef!14] (0,0) rectangle (4.2,1.8);
  \fill[omThm!14] (4.2,0) rectangle (6.4,1.8);
  \draw[thick] (0,0) rectangle (6.4,1.8);
  \draw[thick] (4.2,0) -- (4.2,1.8);
  \node[below, font=\small] at (2.1,0) {$a$};
  \node[below, font=\small] at (5.3,0) {$b$};
  \node[left, font=\small] at (0,0.9) {$k$};
  \node[font=\small] at (2.1,0.9) {$k \times a$};
  \node[font=\small] at (5.3,0.9) {$k \times b$};
\end{tikzpicture}

{\small 分配律的\omterm{def:g6:measure:area}{面积}解释：大长方形
$k \times (a + b)$ 恰为两个较小长方形\omterm{def:g6:measure:area}{面积}之和。}
\end{omfigure}

\begin{example}[两个方向]\label{ex:g7:priorities:distributivity}
\emph{展开}（从左到右）使心算变得容易：
\[
7 \times 103 = 7 \times (100 + 3) = 700 + 21 = 721,
\qquad
6 \times 98 = 6 \times (100 - 2) = 600 - 12 = 588 .
\]
\emph{提取公因式}（从右到左）可把公因式集中起来：
\[
23 \times 12 + 23 \times 8 = 23 \times (12 + 8) = 23 \times 20 = 460 .
\]
\end{example}

\begin{example}[把文字译成算式]\label{ex:g7:priorities:words}
“$8$ 与 $5$ 乘 $3$ 之积的和”是 $8 + 5 \times 3 =
23$。“$5$ 乘以 $8$ 与 $3$ 的和”是
$5 \times (8 + 3) = 55$：文字“……的积乘以……
与……的和”给和加上了括号。
\end{example}

\section{练习}

\begin{exercise}[$\star$]\label{exo:g7:priorities:1}
计算，每行只做一步运算：
\[
5 + 3 \times 6, \qquad
(5 + 3) \times 6, \qquad
30 - 12 \div 4, \qquad
(30 - 12) \div 4 .
\]
\end{exercise}

\begin{exercise}[$\star$]\label{exo:g7:priorities:2}
计算：
\[
24 - 6 + 2, \qquad
24 \div 6 \times 2, \qquad
24 - 6 \times 2, \qquad
24 \div (6 \times 2).
\]
\end{exercise}

\begin{exercise}[$\star$]\label{exo:g7:priorities:3}
逐步计算：
\[
7 + 3 \times (10 - 6), \qquad
40 - (2 + 3) \times 4, \qquad
60 \div \bigl(2 \times (7 - 4)\bigr).
\]
\end{exercise}

\begin{exercise}[$\star$]\label{exo:g7:priorities:4}
计算下列\omterm{def:g6:fractions:def}{分数}（注意\omterm{def:g6:fractions:def}{分数}线相当于括号！）：
\[
\frac{18 + 6}{8}, \qquad
\frac{30}{7 - 2}, \qquad
\frac{5 \times 8}{4 + 6} .
\]
\end{exercise}

\begin{exercise}[$\star$]\label{exo:g7:priorities:5}
利用分配律口算：
\[
8 \times 102, \qquad
25 \times 99, \qquad
14 \times 12 + 14 \times 8, \qquad
37 \times 45 - 37 \times 35 .
\]
\end{exercise}

\begin{exercise}[$\star$]\label{exo:g7:priorities:6}
用一行写出算式，必要时加括号，再
计算：“$50$ 与 $6$ 乘 $7$ 之积的差”；
“$4$ 与 $5$ 的和乘以 $10$ 与
$8$ 的差”。
\end{exercise}

\begin{exercise}[$\star$]\label{exo:g7:priorities:7}
阿米尔买了 $3$ 本笔记本，每本 $2.50$ 元，又买了 $2$ 支钢笔，每支 $1.20$ 元。
用一个算式表示他的总花费，然后计算。
\end{exercise}

\begin{exercise}[$\star\star$]\label{exo:g7:priorities:8}
在 $4 + 2 \times 5 - 1$ 中添加括号，使结果分别为：(a) $13$；
(b) $29$；(c) $24$。（三个结果中有一个完全不需要括号——
是哪一个？）
\end{exercise}

\begin{exercise}[$\star\star$]\label{exo:g7:priorities:9}
$12 + 4 \times k$ 与 $16 \times k$ 对每个数 $k$ 是否相同？
代入 $k = 1$ 和 $k = 2$ 检验，并作出结论。哪一种运算顺序上的错误
会使前者变成后者？
\end{exercise}

\begin{exercise}[$\star\star$]\label{exo:g7:priorities:10}
一张电影票售价 $9$ 欧元，而 $n$ 人团体
需付固定费用 $12$ 欧元，另加每人 $6$ 欧元。
\begin{enumerate}
  \item 写出团体票价的表达式，并计算
        $n = 5$ 时的票价。
  \item 从多少人起，团体票比 $n$
        张单独购票便宜？（试取 $n$ 的几个值。）
\end{enumerate}
\end{exercise}

\begin{exercise}[$\star\star\star$]\label{exo:g7:priorities:11}
每个数字 $1$、$2$、$3$、$4$ 恰好用一次，并任意使用
四则运算与括号，写出分别等于 $24$、
$10$ 与 $1$ 的算式。
\end{exercise}

\section{周末问题：仿效古埃及书吏做乘法}

\begin{problem}[{周末问题——口算技巧（源于
分配律），以及使用了
四千年的加倍—折半乘法}]
\label{pb:g7:priorities:1}
早在计算器出现之前，甚至早在学校开始
操练乘法口诀之前，古埃及书吏与俄罗斯农民
都能仅用\emph{加倍、折半和
相加}完成任意两数的相乘。他们的方法看似魔术，其实
正是分配律（\cref{thm:g7:priorities:distributivity}）
全力发挥作用。本题先锻炼你的心算，
再教你这一古老方法，然后揭开外壳，看清
它究竟为何能够成立。

\medskip
\noindent\textbf{第一部分——不计算的艺术。}
\begin{enumerate}
  \item 仿照
        \cref{ex:g7:priorities:distributivity}，心算下列各式：
        $7 \times 102$ 与 $9 \times 98$。
  \item 心算，并提取公\omterm{def:g5:division:multiple}{因数}：
        $17 \times 13 + 17 \times 7$ 与
        $45 \times 101 - 45$。
  \item “乘以 $5$”技巧：要乘以 $5$，先乘以
        $10$ 再折半。用此法计算 $5 \times 87$，并
        说明该技巧为何成立。
  \item 另有两个技巧需要说明并加以运用：乘以 $25$，可先
        乘以 $100$ 再除以 $4$（计算
        $25 \times 32$）；乘以 $99$，可先乘以 $100$
        再减去一次（计算 $99 \times 47$）。
  \item 警示：分配律是把一个\omterm{def:g5:division:multiple}{因数}分配到一个
        \emph{和}上，绝不能分配到积上。比较
        $8 \times (5 + 3)$ 与 $8 \times 5 + 8 \times 3$，再
        比较 $8 \times (5 \times 3)$ 与
        $(8 \times 5) \times (8 \times 3)$：哪一组相等，
        哪一组不相等？
\end{enumerate}

\medskip
\noindent\textbf{第二部分——加倍—折半法。}
下面给出计算 $13 \times 24$ 的步骤。写出两列。左
列：从 $13$ 开始，\emph{折半}，舍去\omterm{def:g3:division:remainder}{余数}，
一直折到 $1$：\ $13,\ 6,\ 3,\ 1$。右列：从 $24$ 开始，
按同样次数\emph{加倍}：\ $24,\ 48,\ 96,\ 192$。划去
\emph{左}列数为\omterm{def:g3:numbers:evenodd}{偶数}的各行，再把保留下来的
右列数相加。
\begin{enumerate}[resume]
  \item 按上述步骤计算 $13 \times 24$：哪一行被划去，
        剩下的和是多少，它是否等于
        $13 \times 24$？
  \item 用此方法计算 $21 \times 17$（将 $21$ 折半，将
        $17$ 加倍），并用普通乘法检验结果。
  \item 用此方法计算 $16 \times 25$。保留下来的行有几行，
        以及左列从 $1$，$2$，$4$，$8$，$16, \dots$ 开始时
        有何特点？
  \item 现在逐步说明\emph{为何}这一方法成立。每一步的等式
        都用分配律说明：
        \[
        13 \times 24 = (12 + 1) \times 24
        = 12 \times 24 + 24 = 6 \times 48 + 24 .
        \]
        结论：左列折半、右列加倍，乘积保持不变——
        但当左列数为\emph{\omterm{def:g3:numbers:evenodd}{奇数}}时，右列数中会分出一份，
        必须记下。
        表中第一行记下的是哪个数？
  \item 沿表格继续记录：
        $6 \times 48 = 3 \times 96$，然后
        $3 \times 96 = 1 \times 192 + 96$，然后
        $1 \times 192 = 192$。把沿途记下的所有数
        收集起来，并说明这些被记下的数为何
        恰好是左列为\omterm{def:g3:numbers:evenodd}{奇数}的各行中的右列数，
        即步骤所保留的那些行。
\end{enumerate}

\medskip
\noindent\textbf{第三部分——揭开外壳：若干
加倍之和。}
\begin{enumerate}[resume]
  \item 在第 6 题中，保留下来的各行分别对应 $1$、$4$
        和 $8$ 次加倍后的 $24$，方法悄然写下了
        $13 \times 24 = (1 + 4 + 8) \times 24$。验证
        $1 + 4 + 8 = 13$，并在 $1$、$2$、
        $4$、$8$、$16$ 中找出第 7 题为 $21$ 保留的加\omterm{def:g5:division:multiple}{倍数}。
  \item 将 $25$ 与 $42$ 写成来自加\omterm{def:g5:division:multiple}{倍数}列的若干数之和，
        该数列为 $1, 2, 4, 8, 16, 32$，每个数至多用一次。
  \item 说明为何\emph{每一个}整数都能这样表示：
        取不超过它的最大加\omterm{def:g5:division:multiple}{倍数}，减去后重复此过程——
        为何此过程必定停止？写出 $100$ 为
        这样的和式。
  \item 古埃及书吏正是如此做的，早在四千年前：
        为计算 $19 \times 35$，他们制作了
        $35$ 的加倍表（$1 \to 35$，$2 \to 70$，
        $4 \to 140$，$8 \to 280$，$16 \to 560$），挑选
        加起来恰为 $19$ 的各行，然后相加。
        用他们的方法完成计算。
  \item 结尾：现代计算机用 $0$ 和
        $1$ 计数，其乘法本质上与书吏相同，即通过
        加倍与加法。计算机计算 $13 \times 24$ 时需要做哪三次
        加法？（每个数为何都能分解成加倍之和，
        以及 $0$ 与 $1$ 与此有何关系，
        是“乘方”一章中的
        \cref{pb:g8:powers:1}。）
\end{enumerate}
\end{problem}
